\section*{Chapter \ref{ch:b2:basins-compared} --- Sedimentary Basins Compared}

\begin{solution}{exo:b2:basins-compared:1}
(a) Intracratonic: slow subsidence in pulses over hundreds of millions of years, without major
faulting; its cause is debated. (b) Rift: mechanical thinning of the crust during stretching,
then thermal contraction as the lithosphere cools. (c) Passive margin: the thermal subsidence
inherited from the rift stage, amplified by the load of a thick sediment wedge. (d) and (e)
Foreland: the lithosphere bent down by the load of the thrust belt, the Alps and the Himalaya.
\end{solution}

\begin{solution}{exo:b2:basins-compared:2}
(a) $Y = 3000\times(3330 - 2300)/(3330 - 1030) = 3000\times1030/2300 = \qty{1343}{m}$.
(b) $1343 + 200 = \qty{1543}{m}$. (c) Without its load, the basement would lie at \qty{1343}{m}:
$1 - 1343/3000 = \qty{55}{\percent}$ of the thickness is made room for by the load.
\end{solution}

\begin{solution}{exo:b2:basins-compared:3}
(a) $\beta = 35/25 = 1.4$. (b) $125/1.4 = \qty{89.3}{km}$. (c) $S_i = 4334\times(1 - 1/1.4) =
\qty{1238}{m}$.
\end{solution}

\begin{solution}{exo:b2:basins-compared:4}
(a) $\tau = (1.25\times10^{5})^2/(\pi^2\times8\times10^{-7}) = \qty{1.979e15}{s}$, that is
$1.979\times10^{15}/3.156\times10^{13} = \qty{62.7}{Myr}$. (b) $62.7\times0.8 = \qty{50.2}{Myr}$.
(c) $62.7\times(100/125)^2 = \qty{40.1}{Myr}$. (d) $1 - e^{-1} = \qty{63}{\percent}$; $1 -
e^{-3} = \qty{95}{\percent}$.
\end{solution}

\begin{solution}{exo:b2:basins-compared:5}
(a) $0.63\,e^{-1000/1960} = 0.378$; $0.63\,e^{-2000/1960} = 0.227$. (b) $T_s = 1000 -
0.63\times1960\times(0.6004 - 0.3605) = \qty{704}{m}$. (c) $m = 2700\times704 + 1030\times296 =
\qty{2.206e6}{kg.m^{-2}}$; $\bar\rho = \qty{2206}{kg.m^{-3}}$.
\end{solution}

\begin{solution}{exo:b2:basins-compared:6}
(a) $z_2' \leftarrow 704 + 1234.8\,(1 - e^{-z_2'/1960})$: 1077, 1226, 1278, 1296, 1301, 1303,
1304; the layer is \qty{1304}{m} thick. (b) $(2700\times704 + 1030\times600)/1304 =
\qty{1932}{kg.m^{-3}}$. (c) By \qty{304}{m}, or \qty{30}{\percent}: the water that compaction
has squeezed out.
\end{solution}

\begin{solution}{exo:b2:basins-compared:7}
(a) $1 - e^{-100/62.7} = 0.7971$. For $\beta = 1.5$: $3207\times0.4135\times0.7971 =
\qty{1057}{m}$; for $\beta = 2$: $3207\times0.6366\times0.7971 = \qty{1627}{m}$. (b) Yes: 1100 to
\qty{1400}{m} lies between. (c) $\frac{\beta}{\pi}\sin\frac{\pi}{\beta} = 1250/(3207\times0.7971)
= 0.4890$, so $\beta = 1.6 + 0.1\times(0.4890 - 0.4705)/(0.5204 - 0.4705) = 1.64$, within the
50 to \qty{100}{\percent} of stretching inferred for the basin.
\end{solution}

\begin{solution}{exo:b2:basins-compared:8}
(a) The density of the fill replaces $\rho_w$ in the denominator. (b) The numerator is
$125\,000\times(1/3)\times74.69 = \num{3.112e6}$; $S_i = 3.112\times10^{6}/(3184.4 - 2300) =
\qty{3519}{m}$. (c) Under water, $3.112\times10^{6}/2154.4 = \qty{1445}{m}$: the sediment-filled
basin is 2.44 times deeper.
\end{solution}

\begin{solution}{exo:b2:basins-compared:9}
(a) $2171 - 100\times3330/2300 = 2171 - 145 = \qty{2026}{m}$. (b) $Y$ changes metre for metre
with $W_d$: from 1971 to \qty{2371}{m}. (c) The water depth, $\pm\qty{200}{m}$, against
\qty{145}{m} for the sea level.
\end{solution}

\begin{solution}{exo:b2:basins-compared:10}
(a) $S_i$ has the sign of the bracket in its numerator, since the denominator, $3184.4 - 1030$, is
positive; it is negative, an uplift, when the bracket is. (b) $(\rho_m - \rho_c)t_c/a =
\alpha T_1\rho_m/2$ gives $t_c = a\,\alpha T_1\rho_m/(2(\rho_m - \rho_c)) = 125\times0.04372\times
3330/1060 = \qty{17.17}{km}$. (c) $t_c = 17.17/(1 - 0.04372\times17.17/250) = \qty{17.22}{km}$.
(d) Where the crust is thinner than about \qty{17}{km} before stretching: in oceanic lithosphere,
or in a continental crust already thinned by an earlier rifting.
\end{solution}

\begin{solution}{exo:b2:basins-compared:11}
(a) Before stretching, the column holds the lithosphere of mass $M = \rho_c t_c(1 - \alpha T_1
t_c/(2a)) + \rho_m(a - t_c)(1 - \alpha T_1(a + t_c)/(2a))$; after, $\rho_w S_i + M/\beta +
\rho_m(1 - \alpha T_1)(a - S_i - a/\beta)$. Equating them, $S_i[\rho_m(1 - \alpha T_1) - \rho_w] =
(1 - 1/\beta)[\rho_m(1 - \alpha T_1)a - M]$, and expanding $M$, $\rho_m(1 - \alpha T_1)a - M =
(\rho_m - \rho_c)t_c(1 - \alpha T_1 t_c/(2a)) - \rho_m\alpha T_1 a/2$: the formula. (b) The final profile is $T_1 z/a$; the initial one exceeds it
by $T_1 z(\beta - 1)/a$ above $a/\beta$ and by $T_1(1 - z/a)$ below. The integrals are $T_1
a(\beta - 1)/(2\beta^2)$ and $T_1 a(\beta - 1)^2/(2\beta^2)$, whose sum is $T_1 a(1 - 1/\beta)/2$.
The contraction $\alpha$ times this, loaded by water, gives the expression. (c) $3.28\times10^{-5}
\times3330\times1333\times125\,000\times0.5/(2\times2300) = \qty{1978}{m}$, against $3207\times
2/\pi = \qty{2042}{m}$: the first term overstates the final subsidence by \qty{64}{m},
\qty{3}{\percent}.
\end{solution}

\begin{solution}{exo:b2:basins-compared:12}
(a) $3207\times0.6366\times(1 - e^{-150/62.7}) = 2042\times0.9086 = \qty{1855}{m}$. (b) $1855\times
2300/(3330 - 2400) = \qty{4588}{m}$. (c) 6.9 to \qty{13.8}{km}. (d) The model leaves out several
things that could deepen the basin: a mantle lithosphere stretched more than the crust, a regional
(flexural) response to the load of the sediment, a transient uplift of the margin during rifting
followed by sinking, or a supply of sediment so large that the basin was overfilled and loaded
beyond sea level on the shelf.
\end{solution}

\begin{solution}{pb:b2:basins-compared:1}
\textbf{1.} $0.49\,e^{-4420/3704} = 0.149$ at the top; $0.49\,e^{-6030/3704} = 0.096$ at the base.
\textbf{2.} $T_s = 1610 - 0.49\times3704\times(0.3032 - 0.1963) = \qty{1416}{m}$.
\textbf{3.} Shale: $1110 - 0.63\times1960\times(e^{-3310/1960} - e^{-4420/1960}) = 1110 - 1234.8
\times(0.1847 - 0.1049) = \qty{1011}{m}$.
\textbf{4.} Shaly sand: $2680\times1004 + 1030\times696 = \qty{3.408e6}{kg.m^{-2}}$; chalk:
$2710\times1409 + 1030\times201 = \qty{4.025e6}{kg.m^{-2}}$; shale: $2700\times1011 + 1030\times99 =
\qty{2.832e6}{kg.m^{-2}}$; sand: $2650\times1416 + 1030\times194 = \qty{3.952e6}{kg.m^{-2}}$. Total
\qty{1.4217e7}{kg.m^{-2}}, $\bar\rho_s = 14\,216\,910/6030 = \qty{2358}{kg.m^{-3}}$.
\textbf{5.} $z_2' \leftarrow 1416 + 1815\,(1 - e^{-z_2'/3704})$: 1993, 2171, 2221, 2235, 2238, 2239,
2239: \qty{2239}{m}.
\textbf{6.} $(2650\times1416 + 1030\times823)/2239 = \qty{2055}{kg.m^{-3}}$.
\textbf{7.} $Y = 2239\times(3330 - 2055)/2300 + 200 = 1241 + 200 = \qty{1441}{m}$.
\textbf{8.} $Y = 6030\times(3330 - 2358)/2300 + 100 = 2548 + 100 = \qty{2648}{m}$. The units are
measured at their present depths: they are already as compacted as they are today.
\textbf{9.} $(2171 - 1441)/50 = \qty{14.6}{m}$ per million years from 150 to \qty{100}{Ma};
$(2453 - 2171)/40 = \qty{7.05}{m}$ from 100 to \qty{60}{Ma}; $(2648 - 2453)/60 = \qty{3.25}{m}$
since: the sinking slows down.
\textbf{10.} $6030 + 100 = \qty{6130}{m}$, 2.31 times the tectonic subsidence.
\textbf{11.} $(6130 - 2648)/6130 = \qty{57}{\percent}$.
\textbf{12.} $2171 - 100\times3330/2300 = \qty{2026}{m}$.
\textbf{13.} The point at the end of rifting, \qty{150}{Ma}: $S_i = \qty{1441}{m}$, if no thermal
subsidence took place during the \qty{15}{Myr} of rifting.
\textbf{14.} $\beta = 1/(1 - 1441/4334) = 1/0.6675 = 1.50$.
\textbf{15.} $35/1.50 = \qty{23.3}{km}$ of crust and $125/1.50 = \qty{83.3}{km}$ of lithosphere.
\textbf{16.} $(1.50/\pi)\sin(\pi/1.50) = 0.4775\times0.8660 = 0.4135$; $3207\times0.4135 =
\qty{1326}{m}$.
\textbf{17.} $S_i = 4334/3 = \qty{1445}{m}$. At \qty{100}{Ma}, $t = \qty{50}{Myr}$: $1326\times(1 -
e^{-50/62.7}) = \qty{729}{m}$, $Y = \qty{2173}{m}$; at \qty{60}{Ma}: $1326\times0.7620 =
\qty{1010}{m}$, $Y = \qty{2455}{m}$; today: $1326\times0.9086 = \qty{1205}{m}$, $Y = \qty{2650}{m}$.
\textbf{18.} The model is \qty{2}{m} deeper at each point: 2173 against 2171, 2455 against 2453,
2650 against 2648.
\textbf{19.} Palaeo water depths known only to within a hundred metres or more; past sea levels
different from today's; a rifting that lasted \qty{15}{Myr}, during which part of the heat was
already lost; flexure rather than local isostasy; a mantle stretched more or less than the crust;
a crust not \qty{35}{km} thick before rifting.
\textbf{20.} $1326\times(e^{-150/62.7} - e^{-200/62.7}) = \qty{67}{m}$.
\textbf{21.} $66.6\times2300/(3330 - 2400) = \qty{165}{m}$.
\textbf{22.} $1 - e^{-150/62.7} = \qty{91}{\percent}$.
\textbf{23.} $\tau\ln2 = \qty{43.5}{Myr}$ after the end of rifting: at \qty{106.5}{Ma}.
\textbf{24.} $\beta = \textbf{1.50}$: the crust was stretched to one and a half times its width.
\end{solution}
